%----------------------------------------------------------------------------------------
%	Paquets et macros du contenu du rapport TravelHub.
%	Chargé après includes.tex : la mise en page du modèle reste inchangée.
%----------------------------------------------------------------------------------------
\usepackage{listings}
\usepackage{booktabs}
% Colonnes X alignées à gauche : pas d'espaces étirés dans les cellules étroites.
\renewcommand{\tabularxcolumn}[1]{>{\raggedright\arraybackslash}p{#1}}
\usepackage{tikz}
% babel : neutralise les caractères actifs du français (: ; ! ?) dans les figures TikZ.
\usetikzlibrary{positioning,arrows.meta,shapes.multipart,babel}

% Code en ligne : identifiants, chemins, noms de contraintes (_ { } sans échappement).
% Coupure de ligne permise après « _ », « . » et « / » : les longs identifiants
% (trips_participant_single_owner, apps.core.exceptions…) ne débordent pas.
\ExplSyntaxOn
\tl_new:N \l__th_code_tl
\NewDocumentCommand{\code}{m}{
  \tl_set:Ne \l__th_code_tl { \tl_to_str:n {#1} }
  \regex_replace_all:nnN { ([\_\./]) } { \1 \c{allowbreak} } \l__th_code_tl
  \regex_replace_once:nnN { \c{allowbreak} \Z } { } \l__th_code_tl
  \texttt{ \l__th_code_tl }
}
\ExplSyntaxOff
% Les mots en chasse fixe sont larges : souplesse supplémentaire pour les seuls
% paragraphes qui déborderaient sinon.
\setlength{\emergencystretch}{3em}

% Capture d'écran : l'image figures/<fichier> si elle existe, sinon un cadre qui
% indique le fichier attendu (le document compile avant l'ajout des captures).
\newcommand{\capture}[3]{%
  \begin{figure}[H]
    \centering
    \IfFileExists{figures/#1}{%
      \includegraphics[width=\linewidth,height=0.42\textheight,keepaspectratio]{figures/#1}%
    }{%
      \fbox{\parbox[c][4cm][c]{0.9\linewidth}{\centering\small
        Capture à insérer :\\[2pt] \texttt{figures/#1}}}%
    }
    \caption{#2}\label{#3}
  \end{figure}}

% Extraits de code
\definecolor{codegris}{gray}{0.45}
\definecolor{codebleu}{rgb}{0.10,0.25,0.55}
\definecolor{codevert}{rgb}{0.10,0.45,0.20}
\definecolor{codefond}{gray}{0.97}

\renewcommand{\lstlistingname}{Extrait}
% Un extrait n'est jamais coupé entre deux pages (ils font tous moins d'une page).
\BeforeBeginEnvironment{lstlisting}{\par\noindent\begin{minipage}{\linewidth}}
\AfterEndEnvironment{lstlisting}{\end{minipage}\par\smallskip}
\renewcommand{\lstlistlistingname}{Liste des extraits}

\lstdefinelanguage{json}{
  morestring=[b]",
  sensitive=true,
}
\lstdefinelanguage{graphql}{
  morekeywords={query,mutation,fragment,on,true,false,null},
  sensitive=true,
  morecomment=[l]{\#},
  morestring=[b]",
}

\lstset{
  basicstyle=\ttfamily\scriptsize,
  keywordstyle=\color{codebleu},
  commentstyle=\color{codegris}\itshape,
  stringstyle=\color{codevert},
  backgroundcolor=\color{codefond},
  frame=single,
  rulecolor=\color{black!25},
  framesep=3pt,
  xleftmargin=4pt,
  xrightmargin=4pt,
  columns=fullflexible,
  keepspaces=true,
  showstringspaces=false,
  breaklines=true,
  breakatwhitespace=true,
  tabsize=4,
  captionpos=t,
  aboveskip=8pt,
  belowskip=6pt,
  upquote=false,
  extendedchars=true,
  literate=
    {é}{{\'e}}1 {è}{{\`e}}1 {ê}{{\^e}}1 {ë}{{\"e}}1
    {à}{{\`a}}1 {â}{{\^a}}1 {ç}{{\c{c}}}1 {î}{{\^i}}1 {ï}{{\"i}}1
    {ô}{{\^o}}1 {ù}{{\`u}}1 {û}{{\^u}}1
    {É}{{\'E}}1 {À}{{\`A}}1 {Ç}{{\c{C}}}1
    {«}{{\guillemotleft}}1 {»}{{\guillemotright}}1
    {’}{{'}}1 {—}{{---}}1 {…}{{...}}1 {→}{{$\rightarrow$}}1,
}

% Diagrammes
\tikzset{
  entite/.style={
    draw, thick, rectangle split, rectangle split parts=2,
    rectangle split part fill={black!10,white},
    rectangle split part align={center,left},
    text width=4.3cm, inner sep=3pt, font=\scriptsize\ttfamily},
  rel/.style={thick, -{Latex[length=2mm]}},
  card/.style={font=\scriptsize, inner sep=1.5pt},
  etape/.style={draw, rounded corners=2pt, fill=black!4, text width=6.2cm,
    align=center, font=\small, minimum height=0.75cm},
  note/.style={font=\scriptsize, text width=6.4cm, align=left},
}
